\chapter{Names, Types, and Diagnostics}
\label{ch:names-types-diagnostics}

A program that parses is not yet a program that makes sense.
The text \code{x + 1} parses perfectly and means nothing if \code{x} is never defined.
The text \code{if 3 then print 1 else print 0} parses too, yet it asks a number to act as a condition.
This chapter adds the static checks that catch such mistakes before execution.
It resolves every name to its declaration, checks that every operation receives the right kind of value, and reports failures as \term{diagnostics} that point at the offending code.

\section{What names mean}
\label{sec:ch05-what-names-mean}

Consider the Mini fragment \code{x = 1; let x = 2 in print x; print x}, with an existing outer integer binding for \code{x}.
The two prints display different numbers even though both mention \code{x}.
The first print sits inside the \code{let} and sees the inner binding, so it prints $2$.
The second print sits outside and sees the outer binding, so it prints $1$.
\term{Name resolution} assigns each variable use to exactly one declaration using the nesting of scopes.
A \term{scope} is a region of the program in which a set of declarations is visible.
Inner scopes may \term{shadow} outer declarations by reusing a name, as the example shows.
Uses in the inner region refer to the inner declaration, and uses elsewhere are unaffected.
Each use must be resolved before its declaration's type is consulted; resolution and checking can share one traversal.
The fragment \code{print y} with no declaration of \code{y} fails resolution, so asking for the type of \code{y} is meaningless.
Mini types are deliberately tiny: every expression is either an integer or a boolean.
Arithmetic operators take integers and produce integers.
Comparisons take integers and produce booleans.
Conditions in \code{if} and \code{while} must be booleans, and both branches of a conditional must agree where the language requires a single result kind.
These rules fit in one page yet already catch the classic errors: adding a boolean, comparing with a missing variable, and branching on a number.
Table~\ref{tab:ch05-diagnostics} lists the diagnostic kinds Mini reports with an example trigger and the message shape for each.
\begin{table}
\caption{Diagnostic kinds in Mini with a triggering example and the message each one produces.}
\label{tab:ch05-diagnostics}
\centering
\begin{tabular}{lll}
\toprule
Kind & Triggering example & Message shape \\
\midrule
Undefined variable & \code{print y} & undefined name \code{y} at line $L$ \\
Redefinition & \code{let x = 1 in let x = 2} & allowed shadowing, no error \\
Type mismatch & \code{1 + (2 > 3)} & operator needs integers, found boolean \\
Bad condition & \code{if 3 then print 1} & condition must be boolean, found integer \\
Arity error & \code{print} with no operand & \code{print} needs one operand \\
\bottomrule
\end{tabular}
\end{table}
The table distinguishes hard errors from legal shadowing on purpose.
Shadowing is a feature with clear rules, not an accident, so the second row reports no error.
Every other row is a hard error that stops compilation of the enclosing unit.
The message shape always includes the construct, the values or kinds involved, and the position.
Positions come from the spans attached during lexing and parsing, which finally pay for themselves here.
\section{Scopes as a stack}
\label{sec:ch05-scopes-stack}

Resolution walks the tree carrying a stack of scopes from the outermost region to the current one.
Entering a \code{let} pushes its binding, and leaving the body pops it.
Looking up a name scans the stack from the top down and stops at the first match.
This single discipline handles nesting of any depth with no special cases.
Figure~\ref{fig:ch05-scope-chain} shows the stack while resolving the inner print of the shadowing example.
\begin{figure}
\centering
\begin{tikzpicture}[node distance=1.4cm]
\node[basic block] (global) {global: x maps to int};
\node[basic block, below of=global] (letb) {let body: x maps to int};
\node[basic block, below of=letb] (use) {use of x};
\draw[cfg edge] (use) -- node[right, font=\sffamily\scriptsize] {lookup} (letb);
\draw[cfg edge] (letb) -- node[right, font=\sffamily\scriptsize] {outer} (global);
\end{tikzpicture}
\caption{Resolving the inner use of \code{x} finds the nearest binding first and consults the outer scope only on a miss.}
\label{fig:ch05-scope-chain}
\end{figure}
The use points at the \code{let} binding because that frame sits on top of the stack.
The global frame remains reachable for other names but is not consulted for \code{x}.
When the walk leaves the \code{let} body, its frame is popped and later uses of \code{x} find the global binding again.
Types are checked during the same walk with one synthesized kind per expression node.
A number literal has kind integer and a comparison has kind boolean regardless of context.
An arithmetic node checks that both children are integers and reports a mismatch naming the offending child kind.
A conditional checks that its test is boolean and then checks both branches under the same scope stack.
This interleaving is safe because resolution of a subtree finishes before its kind is ever needed.

\section{The checker}
\label{sec:ch05-checker}

Algorithm~\ref{alg:ch05-check} resolves names and checks types in one pass over the tree.
\begin{algorithm}
\caption{Single-pass name resolution and type checking for Mini.}
\label{alg:ch05-check}
\begin{pseudocode}
\Require syntax tree node $N$, scope stack $S$
\Ensure kind of $N$ with all names resolved, or a diagnostic
\Function{Check}{$N, S$}
\If{$N$ is a number} \State \Return integer
\ElsIf{$N$ is a variable} \State resolve it in $S$ or report an undefined name; \Return its kind
\ElsIf{$N$ is a let binding} \State check the value, push the name with its kind, check the body, then pop
\ElsIf{$N$ is an assignment} \State resolve the target; check that the value has its declared kind; \Return unit
\ElsIf{$N$ is an arithmetic node} \State check both children, require integers, and \Return integer
\ElsIf{$N$ is a comparison} \State check both children, require integers, and \Return boolean
\ElsIf{$N$ is a conditional or loop} \State require a boolean test and check the branches or body
\Else \State check each child in order and combine the resulting kinds
\EndIf
\EndFunction
\end{pseudocode}
\end{algorithm}
The algorithm assumes the tree is well formed and that each scope frame maps names to their declared kinds.
It guarantees that any name it accepts refers to exactly one declaration and that any kind it returns follows from the rules.
The \code{let} case is the only one that changes the stack, and it always pops what it pushed, even when checking the body reports an error.
This balance keeps later lookups correct after an error and allows the checker to continue with sibling subtrees.
Continuing after an error needs care to avoid cascades, where one mistake produces ten messages.
The standard remedy is to give the erroneous node a placeholder kind and suppress further complaints that depend only on that placeholder.
The programmer then sees the root cause first and the noise never appears.
The checker records one diagnostic per independent fault: an undefined name, a kind mismatch with both kinds named, or a bad condition with the actual kind found.
Each diagnostic carries the span of its node so the renderer can underline the exact text.

\section{Worked example}
\label{sec:ch05-worked-example}

We now check a program with shadowing and a planted type error by hand.
\begin{irlisting}
x = 1;
let x = 2 in print x + (x > 0);
print x
\end{irlisting}
The checker starts with the global scope containing \code{x} of kind integer from the first assignment.
It enters the \code{let}, checks the value \code{2} as integer, and pushes a second binding of \code{x} to integer.
The stack now holds $2$ frames and $2$ bindings for \code{x}.
Inside the body it checks the print operand \code{x + (x > 0)}.
The left child \code{x} resolves to the inner binding with kind integer.
The comparison \code{x > 0} checks both sides as integers and yields kind boolean.
The addition then sees an integer on the left and a boolean on the right and reports a type mismatch.
The diagnostic names the operator, the two kinds found, and the span of the addition.
The checker assigns the addition a placeholder kind and continues without further messages for that subtree.
After the body it pops the inner frame, leaving $1$ frame.
The final print resolves \code{x} to the global binding of kind integer and checks cleanly.
The run therefore produces exactly $1$ diagnostic from $3$ statements, $2$ pushes and $1$ pop aside, with $2$ live bindings at the deepest point.
The numbers confirm the mechanism: the stack depth tracks scope nesting, and the single error is reported once at its true location.

\section{Implementation notes}
\label{sec:ch05-implementation}

Scope frames are usually hash tables mapping names to declaration records.
Pushing a frame allocates a table, and popping it discards the table, so shadowing needs no save-and-restore bookkeeping.
An alternative keeps one global table with a per-name stack of declarations, which shares memory across scopes at the cost of trickier cleanup.
A per-name stack gives expected constant-time lookup with hashing; separate scope tables may require probing every enclosing scope.
Both preserve the nearest-binding rule of Figure~\ref{fig:ch05-scope-chain}.
Kinds can be small integers rather than heap objects, since Mini has only two of them.
Real languages need structured kinds with variables, constructors, and sharing, plus unification or constraint solving.
The single-pass shape of Algorithm~\ref{alg:ch05-check} still survives there: resolution first within each node visit, then kind combination.
Diagnostics deserve a dedicated rendering path with severity levels for errors, warnings, and notes.
Errors stop code generation, warnings do not, and notes attach context to a primary message.
Capping the number of reported errors keeps a broken file readable; compilers commonly stop after a few dozen and invite a recompile.
The running time is $O(V \cdot L)$, where $V$ is the number of tree nodes visited and $L$ is the worst-case cost of one scope lookup.
With $D$ separate hash-table frames, lookup can take expected $O(D)$ time even when each individual table probe is expected constant time.
Memory use is $O(D + B)$, where $D$ is the maximum scope nesting depth and $B$ is the number of simultaneously live bindings.
Production front ends add recovery rules so that parsing continues after a missing semicolon and checking continues after a bad type.
The guiding rule is that every diagnostic must suggest a fix the programmer can apply without understanding compiler internals.
\begin{ccwarning}[No cascades]
After one fault, suppress every follow-on complaint that could be mere fallout.
Ten messages for one missing declaration teach the programmer to ignore diagnostics, which defeats the whole chapter.
\end{ccwarning}

\section{Exercises}
\label{sec:ch05-exercises}

\begin{enumerate}
\item \label{ex:ch05-resolve-shadow} Resolve every name in \code{a = 1; let a = 2 in let b = a in print a + b; print a} by hand. Give the binding each use refers to, the stack depth at each point, and the final output. Confirm the maximum depth is $3$ frames.
\item \label{ex:ch05-find-mismatch} Find all static faults in \code{if x + 1 then print x} where \code{x} is undeclared. List each diagnostic in the message shape of Table~\ref{tab:ch05-diagnostics} with kinds and positions. Explain the order your checker reports them and which placeholder kinds it uses.
\item \label{ex:ch05-warn-design} Design a warning for an unused \code{let} binding that never triggers on shadowing uses. State the analysis your warning needs beyond Algorithm~\ref{alg:ch05-check}, its running time in terms of $V$ and $B$, and one program where your warning would be wrong.
\end{enumerate}

\section{Further reading}
\label{sec:ch05-further-reading}

Appel covers name binding and type checking as the semantic phase that turns parsed trees into checked programs, which is the direct model for Algorithm~\ref{alg:ch05-check}, see~\cite{appel-1998-modern}.
The C language definition shows how a production standard specifies scopes, linkage, and type constraints with exact rules for shadowing and conversions, see~\cite{wg14-2011-c11}.
Morgan discusses how front-end checks fit into a compiler that must keep compiling large code bases quickly and noisily only when it matters, see~\cite{morgan-1998-optimizing}.
Safonov treats precise diagnostics as part of compiler trustworthiness, since a checker the programmer distrusts will be worked around rather than obeyed, see~\cite{safonov-2010-trustworthy}.
